% Part: second-order-logic
% Chapter: syntax-and-semantics
% Section: introduction

\documentclass[../../../include/open-logic-section]{subfiles}

\begin{document}

\olfileid[pt]{sol}{syn}{int}

\olsection{Introdução}

Na lógica de primeira ordem, combinamos os símbolos não lógicos de uma dada
linguagem, isto é, suas !!{constant}s, !!{function}s e !!{predicate}s, com
os símbolos lógicos para expressar coisas sobre !!{structure}s de primeira
ordem. Isso é feito por meio da noção de satisfação, que relaciona
!!a{structure}~$\Struct{M}$, uma atribuição de variáveis~$s$ e
!!a{formula}~$!A$: $\Sat{M}{!A}[s]$ vale sse é verdadeiro aquilo que $!A$
expressa quando suas !!{constant}s, !!{function}s e !!{predicate}s são
interpretados como determina $\Struct{M}$, e suas variáveis livres são
interpretadas como determina~$s$. A interpretação da !!{identity}~$\eq$
está incorporada à definição de $\Sat{M}{!A}[s]$, assim como a interpretação
de~$\lforall$ e $\lexists$. A primeira é sempre interpretada como a relação
de identidade no !!{domain}~$\Domain{M}$ da estrutura, e os quantificadores
são sempre interpretados como variando sobre todo o !!{domain}. Mas,
fundamentalmente, a quantificação só é permitida sobre elementos do
!!{domain}; assim, somente !!{variable}s de objeto podem seguir um
quantificador.

Na lógica de segunda ordem, tanto a linguagem quanto a definição de
satisfação são estendidas para incluir variáveis livres e ligadas de função
e de predicado, bem como a quantificação sobre elas. Essas variáveis se
relacionam com !!{function}s e !!{predicate}s do mesmo modo que as variáveis
de objeto se relacionam com !!{constant}s. Elas desempenham o mesmo papel
na formação de termos e !!{formula}s da lógica de segunda ordem, e a
quantificação sobre elas é tratada de maneira semelhante. Na semântica
\emph{padrão}, os quantificadores de segunda ordem variam sobre todos os
objetos possíveis do tipo apropriado (funções de $n$ argumentos de
$\Domain{M}$ para~$\Domain{M}$, no caso de variáveis de função, e relações
de $n$ argumentos, no caso de variáveis de predicado). Por exemplo, embora
$\lforall[\Obj{v_0}][(\Obj{P^1_0}(\Obj{v_0}) \lor \lnot
  \Obj{P^1_0}(\Obj{v_0}))]$ seja uma fórmula tanto na lógica de primeira
quanto na de segunda ordem, nesta última também podemos considerar
$\lforall[\Obj{V^1_0}][\lforall[\Obj{v_0}][(\Obj{V^1_0}(\Obj{v_0}) \lor
    \lnot \Obj{V^1_0}(\Obj{v_0}))]]$ e
$\lexists[\Obj{V^1_0}][\lforall[\Obj{v_0}][(\Obj{V^1_0}(\Obj{v_0}) \lor
    \lnot \Obj{V^1_0}(\Obj{v_0}))]]$. Como não contêm variáveis livres,
elas são !!{sentence}s da lógica de segunda ordem. Aqui, $\Obj{V^1_0}$ é
uma variável de predicado de segunda ordem com $1$ argumento. As
interpretações admissíveis de $\Obj{V^1_0}$ são as mesmas que podemos
atribuir a !!{predicate} de $1$ argumento como $\Obj{P^1_0}$, isto é,
subconjuntos de~$\Domain{M}$. Quantificar sobre elas equivale então a dizer
que $\lforall[\Obj{v_0}][(\Obj{V^1_0}(\Obj v_0) \lor \lnot
  \Obj{V^1_0}(v_0))]$ vale para todas as maneiras de atribuir um subconjunto
de~$\Domain{M}$ como valor de $\Obj{V^1_0}$, ou para pelo menos uma delas.
Como todo conjunto contém ou não contém um dado objeto, ambas são
verdadeiras em qualquer !!{structure}.
\end{document}
